Text-to-motion generation converts natural-language descriptions into human
motion, enabling animation and interactive motion
creation~\cite{zhang2023generating,tevet2022humanmotiondiffusionmodel}.
In an interactive setting, a user may revise the wording of a prompt without
changing the intended action, or edit a specific action attribute. Reliable
control requires distinguishing these two cases. For example, ``walk forward
slowly'' and ``move ahead at a slow pace'' should express similar motions;
replacing ``forward'' with ``backward'' should instead change direction while
retaining the action and speed. We call this requirement \emph{selective
prompt robustness}: invariance to preserving rewrites and sensitivity to
intentional semantic edits.

Standard caption--motion likelihood training provides no explicit supervision
for this relation between prompts. T2M-GPT represents motion with a
vector-quantized variational autoencoder (VQ-VAE) and predicts discrete motion
tokens from CLIP text features~\cite{zhang2023generating,radford2021learning}.
Each caption is supervised by its target motion, but the objective does not
specify which changes in language should preserve that target. SATO addresses
instability under synonyms and small textual perturbations by regularizing
text attention and motion prediction~\cite{chen2024sato}. Extending such
invariance to intentional edits, however, can conflict with the requested
change. The relation between prompts must therefore determine the form of
regularization.

A second challenge is choosing the signal to regularize. Text similarity and
attention alone do not identify which words support a particular generated
motion. We use the likelihood of the target motion's early autoregressive
tokens to compute Grad~$\times$~Input attribution. This yields a
motion-conditioned, model-dependent measure of textual evidence. COMPASS uses
that evidence to weight auxiliary text representations, aligns them for
preserving pairs, and defines attribution-change and locality objectives for
edited and unedited positions in altering pairs. A frozen teacher limits
prediction drift. The framework retains the pretrained motion tokenizer and
generator architecture.

The evaluation distinguishes the proposed objective from the effects that
can be established for its implementations. The T2M-GPT-family source models
use frozen CLIP features and detached attribution gradients. Separate native
text-encoder adaptation updates CLIP for T2M-GPT and T5 for MotionGPT before
retrieval evaluation. All generation tests use classifier-free guidance
(CFG) at 1.3; the FID columns use source encoders before this additional stage,
whereas retrieval and the other columns use adapted encoders. These column
groups describe different generated actions. Across three rewrite levels,
the adapted COMPASS variants have higher retrieval point estimates than their
parents on HumanML3D and Motion-X. Source-encoder FID results instead show
dataset- and backbone-dependent trade-offs. KIT-ML, component comparisons,
and frozen HumanML3D-to-KIT text transfer test the scope of these observations.

To interpret prompt stability, we derive local connections between cosine
alignment, projected Jacobian energy, and token-distribution variation.
The analysis states the extra conditions needed to transfer auxiliary
consistency to the generator. We then examine those conditions with
hidden-state, Jacobian, and gradient-connectivity diagnostics. The detached
configuration has lower measured hidden-state sensitivity, but its attribution
losses have no direct gradient path to the generator. This finding separates
a useful conditional interpretation from a verified causal training mechanism.

Our contributions are threefold:
\begin{itemize}
    \item We formulate selective prompt robustness through preserving and
    altering relations, and introduce COMPASS, which combines motion-conditioned
    attribution with relation-specific objectives and teacher regularization.
    \item We provide caption-locked evaluations across three datasets,
    separating source-encoder quality from adapted-encoder retrieval, and
    report both improvements and configuration-dependent regressions.
    \item We derive conditional local stability bounds and use internal
    diagnostics to test their applicability, identifying the gradient-path
    limitation of the reported detached-attribution configuration.
\end{itemize}